% !TeX TS-program = xelatex
% 虚构演示简历 — 仅用于 hellojobs 测试数据
\documentclass{resume-photo}
\usepackage{xeCJK}
\setCJKmainfont{Noto Serif CJK SC}
\usepackage{fontspec}
\usepackage{enumitem}
\setlist[itemize]{itemsep=0pt, topsep=1pt, parsep=0pt, partopsep=0pt}

\ResumeName{于涛}

\begin{document}

\ResumeContacts{
  138-0016-0016,%
  \ResumeUrl{mailto:yutao@hit.edu.cn}{yutao@hit.edu.cn},%
  \textnormal{哈尔滨工业大学 | 计算机科学与技术 · 硕士 | 2022—2025}%
}

\ResumeTitle

\noindent 数据库与存储方向应用研发，熟悉 MySQL 内核参数调优、Binlog 解析与分库分表中间件；支撑订单库峰值 QPS 8k+。注重可观测性与文档沉淀，习惯在评审阶段拆解风险；参与线上复盘与跨团队联调，英语 CET-6，可阅读官方文档与 RFC（演示数据）。

\section{教育背景}
\ResumeItem
{哈尔滨工业大学}
[\textnormal{计算学部 | 计算机科学与技术 | 硕士}]
[2022.09—2025.06]

\begin{itemize}
  \item 研究方向：分布式事务与一致性。
\end{itemize}


\ResumeItem
{企业开放日与实训}
[\textnormal{短期实训营（演示）}]
[2023—2024]

\begin{itemize}
  \item 参加头部互联网公司 2 周封闭实训，完成从需求到上线的完整小组项目。
  \item 在实训中负责接口设计与联调，小组答辩成绩排名前 20\%。
  \item 获得实训营「优秀学员」与内推绿色通道（测试数据，勿当真）。
  \item 沉淀实训笔记 1.5 万字，含架构图与排错记录。
\end{itemize}



\section{专业技能}
\begin{itemize}
  \item 深入理解 InnoDB、索引、事务隔离级别与锁等待分析。
  \item 熟悉 ShardingSphere、Canal 与数据同步延迟监控。
  \item 了解 RocksDB 与 LSM-Tree 基础读写放大问题。
  \item 熟悉 Linux 日常运维与 Shell，具备日志检索与 CPU/内存突增初步定位能力。
  \item 掌握 Git / CI 与 Code Review，了解 Docker 与 Kubernetes 基本编排与滚动发布。
  \item 了解 OAuth2 / JWT、接口幂等与 REST 规范，具备单元测试与集成测试习惯（JUnit/pytest 等）。
  \item 能阅读英文技术文档，具备跨团队沟通与需求澄清能力（测试简历）。
\end{itemize}

\section{实习经历}
\ResumeItem
{京东}
[\textnormal{数据库中间件实习生}]
[2024.06—2024.12]

\begin{itemize}
  \item \textbf{职责}: 分片键迁移方案评估与双写校验工具。
  \item \textbf{成果}: 灰度窗口内数据不一致行数 <0.001\%。
\end{itemize}


\ResumeItem
{中国工商银行 | 软件开发中心}
[\textnormal{金科研发实习生}]
[2024.01—2024.06]

\begin{itemize}
  \item \textbf{合规}: 在监管要求下完成日志脱敏与字段分级打标方案落地。
  \item \textbf{批处理}: 参与日终对账任务优化，将窗口期从 4.5h 压缩到 3.1h。
  \item \textbf{数据库}: 协助 DBA 分析慢 SQL，增加覆盖索引与分页改写，平均耗时下降 62\%。
  \item \textbf{安全}: 参与渗透测试整改，修复 XSS 与越权访问类问题 6 处。
  \item \textbf{流程}: 熟悉银行版本发布、双人复核与变更窗口制度。
\end{itemize}



\section{项目经历}
\ResumeItem
{慢查询根因分析工具}
[\textnormal{实验室课题}]
[2024.01—2024.08]

\begin{itemize}
  \item 基于 Performance Schema 聚类，误报率较规则引擎下降 35\%。
\end{itemize}


\ResumeItem
{特征存储与在线推理}
[\textnormal{Feast 类实践（演示）}]
[2023.08—2024.01]

\begin{itemize}
  \item \textbf{一致}: 离线在线特征对齐，监控 PSI 漂移。
  \item \textbf{延迟}: 在线特征查询 P99 <6ms，支撑排序模型实时更新。
  \item \textbf{平台}: 与模型服务 gRPC 联调，完成金丝雀发布。
  \item \textbf{文档}: 输出特征字典与血缘说明，供算法与工程共用。
\end{itemize}



\section{个人荣誉}
\begin{itemize}
  \item 哈工大 优秀团员
  \item 中国数据库技术大会 DTCC 学生票分享
  \item 全国计算机等级考试四级（网络工程师）（演示）
  \item 校级「优秀学生干部 / 优秀团员」或技术社团骨干（综合测评前 15\%）
  \item 开源贡献：向知名项目提交被合并的 PR（文档/测试/feature）（演示）
\end{itemize}

\end{document}
